\scanpage{258}
\mspage{11}
% archive sheet number: 258
difficultés \unsure{accessoires|excessives} ici, je vais supposer que $A$ est la cat. des $k$-modules, $k$ un anneau quelconque.

À tout $\unsure{\zeta|\xi}\in\deleted{\unsure{\mathrm{Ob}}}\mathrm{Hot}$ est associé un ``abélianisé''

\[
\zeta_{\mathrm{ab}}\overset{\mathrm{def}}{=}H_{\unsure{\bullet|\cdot}}(\zeta)\in
\underset{\text{ou }D(\mathbb Z\text{-}\mathrm{Mod})}{\underbrace{D(\mathrm{Ab})}}
\quad\begin{aligned}&\text{catégorie dérivée des}\\&\text{groupes abéliens,}\end{aligned}
\]

objet qui n'a de cohomologie qu'en degrés $\leq 0$. Ceci dit, si $M$ est dans $D(k\text{-}\mathrm{Mod})$, on aura

\[
\left\{\begin{aligned}
\zeta\mathbin{\underset{\mathbb D}{\times}}M&\simeq\zeta_{\mathrm{ab}}\overset{\mathbb L}{\otimes}_{\mathbb Z}M\\
\operatorname{hom}_{\mathbb D}(\zeta,M)&\simeq\mathbb R\operatorname{Hom}_{\mathbb Z}(\deleted{\unsure{\zeta}}\zeta_{\mathrm{ab}},M),
\end{aligned}\right.
\tag*{(15)}
\]

\unsure{formules|formule} qui est en \unsure{droit fil} avec la formule de Künneth pour l'homologie

\[
(\zeta\times y)_{\mathrm{ab}}\simeq\zeta_{\mathrm{ab}}\overset{\mathbb L}{\otimes}_{\mathbb Z}y_{\mathrm{ab}},
\tag*{(16)}
\]

ce qui va donner un aspect familier aux formules (9), qui se réduisent :

\[
\left\{\begin{aligned}
\zeta_{\mathrm{ab}}\overset{\mathbb L}{\otimes}_{\mathbb Z}
(y_{\mathrm{ab}}\overset{\mathbb L}{\otimes}_{\mathbb Z}M)
&\simeq(\deleted{\unsure{y}}\zeta_{\mathrm{ab}}\overset{L}{\otimes}_{\mathbb Z}y_{\mathrm{ab}})\overset{L}{\otimes}_{\mathbb Z}M\\
\mathbb R\operatorname{Hom}_{\mathbb Z}(\zeta_{\mathrm{ab}},R\operatorname{Hom}_{\mathbb Z}(y_{\mathrm{ab}},M))
&\simeq R\operatorname{Hom}_{\mathbb Z}(\zeta_{\mathrm{ab}}\overset{L}{\otimes}_{\mathbb Z}y_{\mathrm{ab}},M).
\end{aligned}\right.
\]

% slanted
\margin{expliciter aussi les formules d'adjonction (7) (p.~4).}

D'autre part, existe-t-il bien un $\operatorname{homext}(M,N)$, si $M,N$ \unsure{ds} $D(k\text{-}\mathrm{Mod})$ ? On pense à $\mathbb R\operatorname{Hom}_k(M,N)$, mais pour en faire un type d'homotopie, on doit ``tronquer en degrés $>0$''. On pourrait donc prendre

\[
\operatorname{homext}(M,N)=\tau_{\unsure{\leq|<}0}\bigl(\mathbb R\operatorname{Hom}_k(M,N)\bigr),
\]
